\slajdid{09_2-s056}
\begin{frame}[label=09_2-s056]
Ukupna energija disipirana na invertoru je $E_{dyn}=E_P+E_N=C_LV_{DD}^2$.\par\bigskip
Ako se ovaj proces dešava sa učestanošću $f=\frac1{T_P+T_N}$\par\medskip
gde je $T_P$ trajanje logičke jedinice na ulazu, radi P tranzistor\par\medskip
i $T_N$ trajanje logičke jedinice na ulazu, radi N tranzistor\par\medskip
dinamička disipacija na invertoru je
\[P_{dyn}=C_LV_{DD}^2f\]
Ovaj rezultat je zanimljiv, pošto pokazuje da će dinamička disipacija biti „značajna“ zbog relativno velike parazitne kapacitivnosti i učestanosti na kojoj invertor radi. Ali iz još jednog razloga. Na primer većina procesora je predviđena da radi i na višim učestanostima nego što je deklarisano, odnosno od one na kojoj rade u sistemima. Poznato je da se mogu „overklokovati“ odnosno naterati da rade i na višim učestanostima.
\end{frame}
